% Methane (CH\textsubscript{4}) is agriculture's dominant greenhouse gas, and managed grassland
% livestock systems are a major, poorly-instrumented source of it. Ecosystem-scale CH\textsubscript{4}
% flux is difficult to predict because it is driven by the continuous interaction of non-stationary
% biophysical drivers with discrete, spike-inducing livestock management events, and existing
% machine-learning work in this space addresses only gap-filling (within-distribution interpolation),
% not forecasting (sequential extrapolation) or the further step of scenario-based "what-if" analysis
% that a genuine agricultural digital twin requires. This dissertation addresses that gap using seven
% years of Eddy Covariance CH\textsubscript{4} flux data from the North Wyke Farm Platform, a managed
% beef-cattle grassland research site, through a three-stage pipeline: gap-filling a continuous
% CH\textsubscript{4} series, forecasting that series recursively over a full annual horizon, and
% projecting its response under synthetic management and climate scenarios. Each stage benchmarks
% statistical baselines, tree ensembles, and newer tabular foundation models (TabPFN, TabICLv2) under
% a common, leakage-safe temporal evaluation protocol, with uncertainty quantification and
% interpretability attached throughout. The improved gap-filling pipeline's best configuration reaches
% hourly $R^2$ of 0.40--0.58 across all three towers (with a further tabular-foundation-model
% challenger reaching as high as 0.68 at individual towers), comfortably clearing the $R^2<0.1$
% ceiling reported in the closest published literature for this ecosystem type; the resulting
% forecasting benchmark, the
% first applied to this ecosystem-scale signal at a managed temperate grassland to the author's
% knowledge, beats a day-of-year climatology baseline by a wide margin (MASE = 0.691); and the
% scenario-projection stage produces distinguishable, interpretable CH\textsubscript{4} responses to
% livestock-density and grazing-timing changes under CMIP6-derived climate trajectories to 2050. A
% single finding recurs independently across all three stages, via genuinely different methods each
% time: livestock density -- and specifically cattle density -- is the dominant driver of
% CH\textsubscript{4} flux at this site. Absolute predictive accuracy nonetheless remains modest by
% the standards of many other applied machine-learning domains, a reflection of the genuinely
% spike-dominated, high-variance nature of the target rather than a deficiency of any single model
% family tested; this dissertation is accordingly positioned as a benchmarked, uncertainty-quantified
% foundation for CH\textsubscript{4} forecasting and scenario projection at a managed UK grassland
% site, not a finished, deployment-ready system.
